% Section 5 -- Attack and defend. Fill the <<...>> placeholders from
% scripts/analyze.py output and the files in results/.
% Needs booktabs, listings, xcolor in the preamble (already added for section 3).

\section*{5. Attack and Defend}

All attacks run inside the same Docker topology. Attacker containers are
added through \texttt{docker-compose.attack.yml} and the fog rate limiter
is switched on through \texttt{docker-compose.mitigation.yml}, so the base
\texttt{docker-compose.yml} and the task~1/2 parameters are untouched. The
legitimate cameras keep their normal 10~s interval throughout. Each phase
was run for \texttt{<<DURATION>>}~s and the collector's
\texttt{all.jsonl} was saved per phase under \texttt{results/}.

\subsection*{5.1 DoS attack}

The attacker reuses the camera image with its \texttt{--dos} flag: it
drops the inter-request sleep, always sends the large image, and ignores
the response (\texttt{timeout=0.2}). One such container, \texttt{cam\_dos},
is placed in \texttt{zone\_a\_net} and floods \texttt{fog\_a}.

Under the attack and with no mitigation, the latency of the legitimate
cameras in zone~A rose from a baseline mean of \texttt{<<baseline mean>>}~ms
to \texttt{<<attack mean>>}~ms (p95 \texttt{<<attack p95>>}~ms), and
\texttt{<<n>>} legitimate requests failed. The flood saturates the
single CPU the fog node is limited to (task~1.4), so genuine requests
queue behind attacker requests.

\subsection*{5.2 Mitigation}

The mitigation is a sliding-window rate limiter in the fog node. For each
request it checks the number of requests seen from the same source in the
last \texttt{RATE\_WINDOW\_S} seconds, and rejects with HTTP~429 above
\texttt{RATE\_MAX}. It is controlled by environment variables and is off
by default.

\begin{lstlisting}[caption={Rate check, applied before any processing}]
if not allowed(client_ip):
    log(..., http_status=429)
    return Response(status_code=429)
\end{lstlisting}

Two design choices matter. First, the key is the source IP, not the
\texttt{client\_id}: an attacker that rotates its ID (5.3) would slip past
an ID-based limit, but still shares one address. Second, the check runs
before the image is processed, so a rejected request costs the fog almost
no CPU -- this is what frees capacity for the legitimate clients. Rejected
requests are logged with status 429, so they can be counted from the logs.

The threshold of \texttt{RATE\_MAX=5} per 10~s sits well above the
legitimate rate: a real camera sends roughly one request every 10~s, so
even all three together stay far below the limit and are never rejected.

\begin{table}[h]
\centering
\begin{tabular}{lrrrr}
\toprule
Phase & legit ok & mean ms & p95 ms & legit failed \\
\midrule
Normal operation        & <<...>> & <<...>> & <<...>> & <<...>> \\
Attack, no mitigation    & <<...>> & <<...>> & <<...>> & <<...>> \\
Attack, with mitigation  & <<...>> & <<...>> & <<...>> & <<...>> \\
\bottomrule
\end{tabular}
\caption{Legitimate-client performance across the three required states
(cameras cam\_a, cam\_b, cam\_c).}
\end{table}

With the limiter on, the attacker's requests are rejected at the door and
the legitimate cameras return to near their baseline latency, while
\texttt{<<n>>} attacker requests were rejected with 429.

\subsection*{5.3 Botnet and identity rotation}

\paragraph{Botnet.} Three flooding cameras in zone~A
(\texttt{cam\_dos}, \texttt{cam\_dos2}, \texttt{cam\_dos3}) attack
\texttt{fog\_a} at once. Because the limiter keys on source IP and each
container has its own address, each attacker is limited independently and
the mitigation still holds: legitimate latency stayed at
\texttt{<<botnet mean>>}~ms.

\paragraph{Rotating ID.} \texttt{cam\_rotate} floods while changing its
\texttt{id} header on every request. This defeats any ID-based rule --
including the collector's \texttt{rate\_per\_client\_id} alert, which no
longer sees a repeated identity -- but the IP-based limiter and the
\texttt{rate\_per\_client\_ip} alert still catch it, because the address
does not change.

\paragraph{Was zone~B affected?} The two zones are separate internal
networks with separate fog nodes, so zone~B cameras never contend with
the zone~A flood at the fog. They do share the cloud node over
\texttt{backbone\_net}. Zone~B legitimate latency went from
\texttt{<<baseline zB>>}~ms to \texttt{<<attack zB>>}~ms during the
attack, which indicates <<whether the shared cloud caused measurable
cross-zone impact -- state the conclusion from the numbers>>.

\subsection*{5.4 ID spoofing}

\texttt{cam\_spoof} joins \texttt{zone\_a\_net} and sends under the
identity \texttt{cam\_a}, at a normal rate rather than flooding, so that
detection cannot rely on volume. The collector's \texttt{id\_ip\_mismatch}
rule fires the first time \texttt{cam\_a} is seen from a second address:

\begin{lstlisting}[caption={Alert from \texttt{results/spoof\_alerts.jsonl}}]
<<paste the real id_ip_mismatch line>>
\end{lstlisting}

The spoofed camera is detected even though its request rate stays within
normal limits, because the rule compares identity against source address
rather than counting requests.
